% oss-plantilla: oss-2026.11
% Preámbulo corporativo fijo. No editar en sd-NN.tex: exportar_latex.py verificar lo compara byte a byte.
% Para cambiar el formato, editar esta plantilla u oss.sty y subir la versión de la primera línea.
\PassOptionsToPackage{unicode}{hyperref}
\PassOptionsToPackage{hyphens}{url}
\documentclass[11pt]{article}
\usepackage{amsmath,amssymb}
\usepackage{iftex}
\ifPDFTeX
  % Prism y otros editores en la nube: pdfLaTeX con Helvetica como familia por defecto.
  \usepackage[T1]{fontenc}
  \usepackage[utf8]{inputenc}
  \usepackage{textcomp}
  \usepackage[scaled=0.92]{helvet}
  \renewcommand{\familydefault}{\sfdefault}
  \DeclareUnicodeCharacter{2264}{\ensuremath{\leq}}
  \DeclareUnicodeCharacter{2265}{\ensuremath{\geq}}
  \DeclareUnicodeCharacter{2260}{\ensuremath{\neq}}
  \DeclareUnicodeCharacter{2248}{\ensuremath{\approx}}
  \DeclareUnicodeCharacter{2190}{\ensuremath{\leftarrow}}
  \DeclareUnicodeCharacter{2192}{\ensuremath{\rightarrow}}
  \DeclareUnicodeCharacter{25CF}{\ensuremath{\bullet}}
  \DeclareUnicodeCharacter{2713}{\ensuremath{\checkmark}}
\else
  % XeLaTeX/LuaLaTeX (compilación oficial local): DejaVu Sans.
  \usepackage{fontspec}
  \defaultfontfeatures{Ligatures=TeX,Scale=MatchLowercase}
  \setmainfont{DejaVu Sans}
  \setsansfont{DejaVu Sans}
  \setmonofont{DejaVu Sans Mono}
\fi
\usepackage[letterpaper,margin=20mm]{geometry}
\usepackage{microtype}
\usepackage{parskip}
\usepackage{xcolor}
\usepackage{longtable,booktabs,array}
\usepackage{calc}
\usepackage{etoolbox}
\makeatletter
\patchcmd\longtable{\par}{\if@noskipsec\mbox{}\fi\par}{}{}
\makeatother
\usepackage{footnotehyper}
\makesavenoteenv{longtable}
\usepackage{graphicx}
\makeatletter
\newsavebox\pandoc@box
\newcommand*\pandocbounded[1]{%
  \sbox\pandoc@box{#1}%
  \Gscale@div\@tempa{\textheight}{\dimexpr\ht\pandoc@box+\dp\pandoc@box\relax}%
  \Gscale@div\@tempb{\linewidth}{\wd\pandoc@box}%
  \ifdim\@tempb\p@<\@tempa\p@\let\@tempa\@tempb\fi
  \ifdim\@tempa\p@<\p@\scalebox{\@tempa}{\usebox\pandoc@box}%
  \else\usebox{\pandoc@box}%
  \fi}
\def\maxwidth{\ifdim\Gin@nat@width>\linewidth\linewidth\else\Gin@nat@width\fi}
\def\maxheight{\ifdim\Gin@nat@height>\textheight\textheight\else\Gin@nat@height\fi}
\makeatother
\setkeys{Gin}{width=\maxwidth,height=\maxheight,keepaspectratio}
\makeatletter
\def\fps@figure{htbp}
\makeatother
\usepackage[normalem]{ulem}
\providecommand{\st}[1]{\sout{#1}}
\usepackage{fancyvrb}
\setlength{\emergencystretch}{3em}
\providecommand{\tightlist}{%
  \setlength{\itemsep}{0pt}\setlength{\parskip}{0pt}}
\providecommand{\passthrough}[1]{#1}
\setcounter{secnumdepth}{-\maxdimen}
\ifPDFTeX
  \usepackage[spanish]{babel}
\else
  \usepackage[bidi=default]{babel}
  \babelprovide[main,import]{spanish}
\fi
\let\LanguageShortHands\languageshorthands
\def\languageshorthands#1{}
\usepackage{bookmark}
\usepackage{xurl}
\urlstyle{same}
\hypersetup{pdflang={es-CL},pdfcreator={Only Simple Solutions - exportar_latex.py}}
\usepackage{oss}
\begin{document}

\ossCover{Plan de riesgos}{Subdocumento 8}
\tableofcontents
\clearpage

\section{Capítulo 8 · Introducción a los Riesgos}\label{capuxedtulo-8-introducciuxf3n-a-los-riesgos}

Este capítulo presenta cómo Only Simple Solutions identifica, califica y trata los riesgos de la solución propuesta a Multitiendas Ancoa S.A. durante los 56 meses del contrato. Se organiza en tres apartados que se leen en orden. El apartado 8.1 define el plan de gestión, con el enfoque, los roles, las escalas de probabilidad e impacto, las zonas de exposición y los umbrales de escalamiento. El apartado 8.2 identifica y analiza los riesgos y reúne 37 riesgos y 8 oportunidades, de los cuales 3 son críticos y 11 son altos. El apartado 8.3 define el plan de acción de cada riesgo y las reservas que lo sostienen.

El capítulo se apoya en otros y les entrega insumos. Los supuestos y restricciones del alcance del Capítulo 3 y las decisiones de arquitectura del Capítulo 4 originan buena parte de los riesgos. Los comités del Capítulo 6 reciben los riesgos según su zona y deciden sobre ellos. El cronograma, la estructura de desglose del trabajo y la ruta crítica del Capítulo 7 determinan qué eventos afectan un hito sin holgura, y las acciones de respuesta se incorporan a ese plan. El plan de calidad del Capítulo 9 verifica que las acciones se cumplan. Los riesgos de adopción de las cinco innovaciones obligatorias no se repiten aquí y se tratan en el Capítulo 13.

El registro completo se entrega en el Formulario T-16, cuya versión resumida se incluye como anexo de este capítulo. El monto de las reservas se informa únicamente en la oferta económica.

\section{8.1 Plan de riesgos}\label{plan-de-riesgos}

Este apartado define las reglas con las que Only Simple Solutions gestionará los riesgos de la solución propuesta para Multitiendas Ancoa S.A. durante los 56 meses del contrato. Establece el enfoque, los roles, las escalas de calificación, las zonas de exposición y los umbrales de escalamiento, de modo que cada riesgo se califique con el mismo criterio. El apartado 8.2 aplica estas reglas a la identificación y al análisis de los riesgos, y el apartado 8.3 define las acciones asociadas a cada uno.

\setlength{\tabcolsep}{3pt}

\subsection{Enfoque de gestión}

La gestión de riesgos sigue la norma ISO 31000 (International Organization for Standardization, 2018) y se opera como un proceso continuo y no como un documento que se redacta una sola vez. Los riesgos se identifican, se analizan, se valoran, se tratan, se monitorean y se comunican de manera cíclica, y cada vuelta del ciclo deja evidencia en el registro. El registro es un instrumento vivo, con un identificador, un dueño, una calificación, una estrategia, un disparador y un estado por cada riesgo. Se revisa en cada sesión del Comité de Proyecto, de modo que ningún riesgo permanece sin revisión más de quince días.

El registro se publica en el espacio colaborativo que el proyecto habilita para el CLIENTE, junto con las actas, los entregables y el registro de cambios. Así, la contraparte de Ancoa consulta en cualquier momento el estado de cada riesgo sin solicitarlo al equipo. Además, cada informe mensual de avance incorpora los riesgos del período, con sus cambios de calificación, los riesgos cerrados y los nuevos, junto con las desviaciones y los compromisos del período siguiente. De esta forma, la información de riesgos llega al patrocinador por el mismo canal que el avance físico y financiero.

\subsection{Roles y responsabilidades}

Cada riesgo tiene un dueño, que es el rol del equipo mejor ubicado para actuar sobre su causa. El dueño propone la calificación, ejecuta o supervisa la acción acordada e informa los cambios de estado antes de cada sesión. El Jefe de Proyecto mantiene el registro, consolida las calificaciones, verifica que cada una se justifique con un dato del Caso o del subdocumento fuente y presenta el resumen al Comité de Proyecto.

Tres instancias de gobierno del Capítulo 6 reciben los riesgos según su naturaleza. El Comité de Proyecto, que sesiona quincenalmente con la contraparte técnica, revisa el registro completo y decide sobre los riesgos de zona alta. El Comité de Arquitectura, que sesiona mensualmente con TI y el encargado de seguridad, revisa los riesgos técnicos. El Comité Ejecutivo, que sesiona mensualmente con el patrocinador y el administrador del contrato, recibe los riesgos críticos y decide la liberación de la reserva de gestión. Desde el mes 13, el Comité de Operación, también mensual, incorpora los riesgos de preparación y de operación. Las contrapartes del CLIENTE validan la calificación de los riesgos que afectan su ámbito, y los dueños de producto confirman el impacto sobre la funcionalidad comprometida.

\subsection{Escalas de probabilidad e impacto}

Los riesgos se califican en cinco niveles, desde muy bajo hasta muy alto. La Tabla 8.1 define la probabilidad y los impactos en plazo, esfuerzo y costo, y la Tabla 8.2 define los impactos en alcance y calidad y en servicio.

Tabla 8.1: Escala de probabilidad e impacto en plazo, esfuerzo y costo. Fuente: Elaboración propia.

\begin{longtable}{p{0.158\linewidth}p{0.158\linewidth}p{0.218\linewidth}p{0.237\linewidth}p{0.158\linewidth}}
\toprule
\textbf{Nivel} & \textbf{Probabilidad} & \textbf{Impacto en plazo} & \textbf{Impacto en esfuerzo} & \textbf{Impacto en costo} \\
\midrule
\endhead
\bottomrule
\endlastfoot
Muy alto (5) & $>$ 70\% & $>$ 6 semanas & $>$ 5\% de las horas del proyecto & $>$ 5\% del contrato \\
Alto (4) & 51 a 70\% & 3 a 6 semanas & 2 a 5\% de las horas & 2 a 5\% del contrato \\
Medio (3) & 31 a 50\% & 1 a 3 semanas & 1 a 2\% de las horas & 1 a 2\% del contrato \\
Bajo (2) & 11 a 30\% & 2 días a 1 semana & 0,5 a 1\% de las horas & 0,5 a 1\% del contrato \\
Muy bajo (1) & 1 a 10\% & Menos de 2 días & Menos de 0,5\% de las horas & Menos de 0,5\% del contrato \\
\end{longtable}

El impacto en costo se expresa siempre como porcentaje del contrato, porque el monto corresponde únicamente al Sobre 3. El impacto en esfuerzo se expresa como porcentaje de las horas de la estimación del apartado 7.2.

Tabla 8.2: Escala de impacto en alcance y calidad, y en servicio. Fuente: Elaboración propia.

\begin{longtable}{p{0.16\linewidth}p{0.34\linewidth}p{0.42\linewidth}}
\toprule
\textbf{Nivel} & \textbf{Impacto en alcance y calidad} & \textbf{Impacto en servicio} \\
\midrule
\endhead
\bottomrule
\endlastfoot
Muy alto (5) & Efecto muy significativo sobre la funcionalidad comprometida. & Incumple una restricción del Caso, la separación de los dos negocios o un umbral con sanción, o pierde datos transaccionales. \\
Alto (4) & Efecto significativo sobre la funcionalidad comprometida. & Interrumpe más de la autonomía de 24 h en tienda o incumple el RTO o el RPO. \\
Medio (3) & Efecto sobre áreas funcionales clave. & Degrada de forma visible un servicio crítico, sin pérdida de datos. \\
Bajo (2) & Efecto menor sobre la funcionalidad general. & Degrada brevemente el servicio, con rodeo disponible. \\
Muy bajo (1) & Efecto menor sobre funciones secundarias. & No tiene efecto para el cliente final. \\
\end{longtable}

Para calcular la exposición, cada nivel se traduce a un valor numérico. La probabilidad toma el punto medio de cada rango, es decir, 0,05, 0,20, 0,40, 0,60 y 0,85. El impacto toma los valores 0,05, 0,10, 0,20, 0,40 y 0,80, que son una decisión del proponente y no provienen de las Bases, ya que ISO 31000 deja los criterios de calificación a cada organización. La progresión del impacto es más pronunciada que la de la probabilidad para que los eventos de consecuencia máxima alcancen la zona crítica con una probabilidad moderada.

Cuatro reglas gobiernan la calificación y evitan que dos personas califiquen el mismo riesgo de manera distinta. Primero, si el evento atrasa un hito contractual sin holgura, esto es, el paso a producción del mes 16, el del mes 21, el retiro de la plataforma de 2011 en el mes 22 o el cierre del mes 56, el impacto en plazo es muy alto sin importar las semanas, porque los eslabones principales de la ruta crítica del apartado 7.3.3 tienen holgura cero y las salidas son no negociables. Segundo, todo evento que borre la separación de datos entre el negocio de retail y la filial emisora de crédito es muy alto en servicio, porque el Caso lo declara línea roja. Tercero, se califica el peor de los impactos entre plazo, esfuerzo, costo, alcance y servicio, y la calificación se justifica con un dato del Caso o del subdocumento fuente. Cuarto, las oportunidades se califican con la misma escala de probabilidad y con impacto positivo, y no se suman a las reservas.

\subsection{Matriz de exposición y zonas}

La exposición de cada riesgo es el producto de su probabilidad por su impacto, ambos en valor numérico, y ubica al riesgo en una de cuatro zonas. La Tabla 8.3 define cada zona y lo que obliga.

Tabla 8.3: Zonas de exposición y regla de acción. Fuente: Elaboración propia.

\begin{longtable}{p{0.12\linewidth}p{0.20\linewidth}p{0.40\linewidth}p{0.20\linewidth}}
\toprule
\textbf{Zona} & \textbf{Exposición (P $\times$ I)} & \textbf{Regla de acción} & \textbf{Reflejo en la oferta} \\
\midrule
\endhead
\bottomrule
\endlastfoot
Crítica & $\geq$ 0,40 & Se cambia el diseño para que el riesgo deje de existir o se escala al Comité Ejecutivo; no se acepta. & Arquitectura, alcance o exclusión explícita. \\
Alta & 0,20 a menos de 0,40 & Se ejecuta una acción de mitigación con costo y plazo dentro del plan. & EDT y cronograma. \\
Media & 0,08 a menos de 0,20 & Se acepta activamente, con reserva de contingencia y disparador definido. & Reserva de contingencia. \\
Baja & $<$ 0,08 & Se acepta pasivamente, con observación y revisión mensual. & Registro, sin costo asociado. \\
\end{longtable}

Con los valores adoptados, un riesgo de probabilidad 0,60 e impacto 0,80 tiene exposición 0,48 y es crítico, mientras que uno de probabilidad 0,40 e impacto 0,40 tiene exposición 0,16 y es medio. La regla del hito contractual aumenta el número de riesgos que caen en la zona crítica, y esa consecuencia es intencional, porque el cronograma de 56 meses no admite desviaciones.

\subsection{Apetito, tolerancia y umbral de escalamiento}

El apetito de riesgo es bajo en lo contractual y regulatorio, esto es, en los hitos, en la separación de los dos negocios y en el cumplimiento ante la CMF, y es moderado en lo técnico mejorable, como el rendimiento y la adopción, donde se acepta ajustar el diseño mientras no cambie un hito. La tolerancia es cero atraso en los hitos de los meses 16, 21 y 56 y cero cruces de datos sin base, y en el resto llega hasta lo que cubra la reserva de contingencia.

El umbral de escalamiento traduce lo anterior en instancias de decisión. Todo riesgo en zona crítica pasa al Comité Ejecutivo, todo riesgo en zona alta pasa al Comité de Proyecto, y los riesgos técnicos se presentan además al Comité de Arquitectura. La liberación de la reserva de gestión corresponde únicamente al Comité Ejecutivo.

\subsection{Ciclo de revisión}

La revisión se organiza en cuatro ritmos. El registro completo se revisa cada quince días en el Comité de Proyecto, y los riesgos técnicos y los críticos se revisan cada mes en los comités de Arquitectura y Ejecutivo. Además, el registro se revisa en cada hito contractual y antes de cada una de las ventanas de congelamiento de publicación, momentos en que el costo de una sorpresa es mayor. El informe mensual de avance resume el resultado de estas revisiones para el CLIENTE.

La calificación se reevalúa completa cuando cambia un supuesto del que depende. Si se modifica el supuesto SUP-26, que fija el inicio del proyecto en enero de 2027, se recalculan los riesgos de plazo frente a las ventanas de congelamiento y a los hitos de los meses 16 y 21. Si se modifica el supuesto SUP-28, que considera 14 tiendas sin enlace de respaldo, se recalculan los riesgos de continuidad de servicio en tienda.

El apartado 8.2 aplica estas escalas y zonas a los riesgos identificados y justifica cada calificación, y el apartado 8.3 define para cada riesgo la acción, el responsable, el plazo y el disparador.

\section{8.2 Identificación y Análisis de Riesgos}\label{identificaciuxf3n-y-anuxe1lisis-de-riesgos}

Este apartado identifica y analiza los riesgos de la solución que Only Simple Solutions propone a Multitiendas Ancoa S.A. Los riesgos provienen de tres fuentes. La primera es el Caso, con sus dolores, sus restricciones y sus fechas. La segunda es la solución descrita en los capítulos 3 y 4, con sus servicios, sus decisiones de arquitectura y su despliegue híbrido. La tercera es el plan del capítulo 7, con el cronograma de 56 meses, los hitos sin holgura y las ventanas de congelamiento. Por esa razón, cada riesgo nace de una causa concreta de esta solución y no de un catálogo genérico. El resultado es un registro de 37 riesgos y 8 oportunidades, cuyo detalle completo se presenta en el anexo del Formulario T-16. El cuerpo del apartado desarrolla los riesgos críticos, los altos y los once que el Caso y las Bases exigen abordar, y deja el análisis cuantitativo y las respuestas para los apartados siguientes.

\subsection{Estructura de desglose de riesgos}

Los 37 riesgos se ordenan en una estructura de desglose de riesgos de cinco ramas, que se muestra en la Figura 8.1. La rama técnica agrupa 8 riesgos de la arquitectura, las plataformas y los datos. La rama organizacional agrupa 5 riesgos de las personas y del cliente. La rama de proyecto agrupa 16 riesgos de plazos, ventanas, supuestos y terceros. La rama de seguridad agrupa 3 riesgos y la rama de operación agrupa 5 riesgos del servicio una vez en producción.

\begin{center}
  \includegraphics[width=\linewidth,height=\textheight,keepaspectratio]{figuras/diag-08-01_rbs.pdf}
\end{center}

Figura 8.1: Estructura de desglose de riesgos, con las cinco ramas y los 37 riesgos. Fuente: Elaboración propia.

La figura muestra que la rama de proyecto concentra casi la mitad de los riesgos, lo que es coherente con un cronograma de 56 meses con hitos contractuales sin holgura y con ventanas de congelamiento que limitan las fechas de despliegue. Nueve de esos 16 riesgos tienen origen externo, ya sea normativo o de terceros. Son R-16 (modalidad fiscal de contingencia sin aprobación de la autoridad tributaria), R-17 (residencia de los datos de recuperación), R-19 (reglamentación de la Ley 21.719), R-20 (hardware, enlaces y obras que provee el cliente), R-27 (cooperación de los seis proveedores actuales), R-28 (retención de grabaciones de repactación), R-29 (procedimiento abierto por el SERNAC), R-32 (ubicación de plataformas y enlaces por tienda) y R-36 (adopción de contratos de estado por transportistas y vendedores de marketplace). Se mantienen dentro de la rama de proyecto y se identifican como externos porque el Caso exige que lo excluido del alcance aparezca en el plan y en el riesgo (Multitiendas Ancoa S.A., 2026, Caso 09, cap. 11, p. 27).

\subsection{Cobertura de los riesgos exigidos}

El Caso espera seis riesgos en el plan, que son la divergencia de saldos en la cartera (R-01), el evento anual con un servicio a medio desplegar (R-02), el cruce indebido de datos entre los dos negocios (R-03), la integración no documentada (R-04), la rotación del 62 \% del personal (R-05) y la comisión del vendedor (R-06) (Multitiendas Ancoa S.A., 2026, Caso 09, cap. 19, p. 46). Además, deben abordarse cinco riesgos adicionales, que son la obsolescencia de los servicios administrados (R-07), el bloqueo por proveedor (R-08), la escalabilidad (R-09), la ciberseguridad (R-10) y las contrapartes del cliente (R-11). El registro cubre los once. Cada riesgo se mira desde la solución, el desarrollo o la implantación, y un mismo riesgo puede tocar más de una mirada. En total, 14 riesgos tocan la solución, 13 el desarrollo y 14 la implantación. La Tabla 8.4 sintetiza los once riesgos exigidos con su rama, su exposición y su zona.

Tabla 8.4: Riesgos exigidos por el Caso y por las Bases. Fuente: Elaboración propia.

\begin{longtable}{p{0.097\linewidth}p{0.407\linewidth}p{0.155\linewidth}p{0.136\linewidth}p{0.136\linewidth}}
\toprule
\textbf{ID} & \textbf{Riesgo} & \textbf{Rama} & \textbf{Exposición} & \textbf{Zona} \\
\midrule
\endhead
\bottomrule
\endlastfoot
R-01 & Divergencia de saldos en la migración de cartera & Técnica & 0,32 & Alta \\
R-02 & Evento anual con servicio a medio desplegar & Proyecto & 0,48 & Crítica \\
R-03 & Cruce indebido de datos entre los dos negocios & Seguridad & 0,16 & Media \\
R-04 & Integración no documentada entre las 14 interfaces & Técnica & 0,12 & Media \\
R-05 & Rotación del 62 \% del personal & Organizacional & 0,17 & Media \\
R-06 & Comisión del vendedor y venta con tarjeta & Organizacional & 0,08 & Media \\
R-07 & Obsolescencia de servicios administrados & Técnica & 0,08 & Media \\
R-08 & Bloqueo por proveedor de nube & Técnica & 0,08 & Media \\
R-09 & Escalabilidad en el evento anual & Técnica & 0,16 & Media \\
R-10 & Ciberseguridad y datos de tarjeta & Seguridad & 0,16 & Media \\
R-11 & Validaciones tardías del cliente & Organizacional & 0,32 & Alta \\
\end{longtable}

La tabla muestra que de los once riesgos exigidos uno es crítico, dos son altos y ocho quedan en zona media. Los apartados siguientes explican por qué cada uno recibió esa zona.

\subsection{Análisis cualitativo}

Cada riesgo se calificó con las escalas y zonas del apartado 8.1, y la Figura 8.2 ubica los 37 en la matriz de exposición según su probabilidad y su impacto.

\begin{center}
  \includegraphics[width=\linewidth,height=\textheight,keepaspectratio]{figuras/diag-08-02_matriz-exposicion.pdf}
\end{center}

Figura 8.2: Matriz de exposición con los 37 riesgos ubicados por probabilidad e impacto. Fuente: Elaboración propia.

La matriz distribuye los riesgos en 3 críticos, 11 altos, 21 medios y 2 bajos. Los tres críticos (R-02, R-12 y R-13) coinciden en probabilidad 0,60 e impacto 0,80, y los tres afectan un hito contractual sin holgura. Siete de los once altos tienen impacto 0,80 con probabilidad 0,40, y los otros cuatro combinan probabilidad 0,60 con impacto 0,40. Cinco riesgos (R-02, R-11, R-23, R-33 y R-37) alcanzan impacto 0,80 por la regla del hito contractual, ya que atrasan eslabones de la ruta crítica del apartado 7.3.3, que tienen holgura cero. Los dos riesgos bajos son R-29 y R-38. La Tabla 8.5 presenta los 14 riesgos críticos y altos, ordenados por exposición.

Tabla 8.5: Riesgos críticos y altos. Fuente: Elaboración propia.

\begin{longtable}{p{0.097\linewidth}p{0.465\linewidth}p{0.116\linewidth}p{0.116\linewidth}p{0.136\linewidth}}
\toprule
\textbf{ID} & \textbf{Riesgo} & \textbf{P} & \textbf{I} & \textbf{Exposición} \\
\midrule
\endhead
\bottomrule
\endlastfoot
R-02 & Evento anual con servicio a medio desplegar & 0,60 & 0,80 & 0,48 \\
R-12 & Marcha blanca dentro de los congelamientos & 0,60 & 0,80 & 0,48 \\
R-13 & Cartera sin migrar al retiro de la plataforma de 2011 & 0,60 & 0,80 & 0,48 \\
R-01 & Divergencia de saldos en la migración de cartera & 0,40 & 0,80 & 0,32 \\
R-11 & Validaciones tardías del cliente & 0,40 & 0,80 & 0,32 \\
R-16 & Modalidad fiscal de contingencia sin aprobar & 0,40 & 0,80 & 0,32 \\
R-23 & Salida de un rol clave del equipo & 0,40 & 0,80 & 0,32 \\
R-25 & Falla posterior al retiro, dentro del congelamiento & 0,40 & 0,80 & 0,32 \\
R-33 & Sobreasignación del equipo en los meses 13 a 15 & 0,40 & 0,80 & 0,32 \\
R-37 & Decisiones del cliente demoradas antes de la línea base & 0,40 & 0,80 & 0,32 \\
R-18 & Carga inicial de existencias con errores & 0,60 & 0,40 & 0,24 \\
R-20 & Hardware y enlaces de tienda no listos para su ola & 0,60 & 0,40 & 0,24 \\
R-24 & Primer peak con defectos latentes & 0,60 & 0,40 & 0,24 \\
R-34 & Meta de error de precio con etiquetas en papel & 0,60 & 0,40 & 0,24 \\
\end{longtable}

La tabla permite ver que los riesgos de mayor exposición se concentran en la migración de la cartera, en las ventanas de congelamiento y en las dependencias del cliente. Los apartados siguientes desarrollan cada uno.

\subsection{Riesgos críticos}

Los tres riesgos críticos tienen exposición 0,48 y, según la regla de acción del apartado 8.1, no se aceptan. Se resuelven cambiando el diseño o el plan, o se escalan al Comité Ejecutivo.

El riesgo R-02 se origina en que la asociación gremial fija la fecha del evento anual con unas seis semanas de aviso (Multitiendas Ancoa S.A., 2026, Caso 09, cap. 13, núm. 13.2, p. 29). Si el evento cae con un servicio a medio desplegar o dentro de una marcha blanca, la venta digital puede caer o un congelamiento desplaza entre 3 y 6 semanas el trabajo de la ruta crítica. Afecta al paso a producción del mes 16 y se observa desde la mirada de implantación. Se calificó con probabilidad 0,60, porque el evento se repite cada año, e impacto 0,80, porque el plan del capítulo 7 no tiene holgura y rige la regla del hito contractual. La exposición resultante es 0,48.

El riesgo R-12 nace de que, con inicio en enero de 2027, la marcha blanca de la Etapa 1 (meses 13 a 15) cae casi entera en los congelamientos, que rigen hasta el 6 de enero y desde la última semana de enero hasta la primera de marzo (Multitiendas Ancoa S.A., 2026, Caso 09, cap. 13, núm. 13.2, p. 29). En ese lapso no pueden corregirse defectos, por lo que el cierre se atrasa y arrastra el paso a producción del mes 16, previsto para abril de 2028. El plan del capítulo 7 habilita la marcha blanca en la ventana posterior al 6 de enero y no admite liberaciones funcionales nuevas durante un congelamiento, por lo que una corrección debe esperar a la ventana siguiente. La mirada es de implantación. Se calificó con probabilidad 0,60, porque el calendario ya coincide con las ventanas, e impacto 0,80, porque afecta un hito sin holgura. La exposición es 0,48.

El riesgo R-13 surge porque la plataforma de crédito de 2011 se retira en el mes 22 (octubre de 2028) y la remediación vence en el mes 24, con fin de soporte en 2029 (Multitiendas Ancoa S.A., 2026, Caso 09, cap. 7, núm. 7.3, p. 16). La ventana de migración mide cerca de 197 días intervenibles entre abril y octubre de 2028, con holgura casi nula, y el tramo complejo de la cartera depende de la Etapa 2, que entra en producción en el mes 21. Si un tramo no alcanza a migrar, el atraso cae en el congelamiento y rebasa el plazo anterior a 2029, con soporte vencido o incumplimiento ante la CMF. Las miradas son de desarrollo e implantación. Se calificó con probabilidad 0,60 e impacto 0,80 por el hito del mes 22, de donde resulta una exposición de 0,48.

\subsection{Riesgos altos}

Los once riesgos altos tienen exposición entre 0,24 y 0,32 y exigen una acción de mitigación con plazo dentro del plan. Se agrupan por tema para facilitar su lectura.

Los tres primeros se relacionan con la cartera y los datos que se migran. El riesgo R-01 se origina en que la cartera viva con cuotas, repactaciones y cobranzas migra por tramos, y el tramo difícil lo hace en la Etapa 2 con una sola marcha blanca de dos meses. Una divergencia de saldos incumple la restricción 8 y expone ante la CMF (Multitiendas Ancoa S.A., 2026, Caso 09, cap. 10, restricción 8, p. 26). Tiene probabilidad 0,40 e impacto 0,80 en servicio, con exposición 0,32.

El riesgo R-18 parte de que el registro de inventario tiene 12,4 \% de discrepancia en el conteo cíclico (Multitiendas Ancoa S.A., 2026, Caso 09, cap. 7, núm. 7.1, p. 15) y de que 24 instalaciones no cierran para el corte (Multitiendas Ancoa S.A., 2026, Caso 09, cap. 19, p. 46). La carga inicial del Servicio de existencias hereda esos errores y el disponible publicado no es confiable en el mes 16. Tiene probabilidad 0,60 e impacto 0,40 en alcance y calidad, con exposición 0,24.

El riesgo R-25 se produce porque el retiro de la plataforma de 2011 ocurre en octubre de 2028, antes del congelamiento de noviembre a enero. Si algo falla después, la corrección o el retorno quedan bloqueados hasta el 7 de enero de 2029, con incumplimiento ante la CMF. La falla requiere que antes falle otro elemento, por lo que la probabilidad es 0,40. Con impacto 0,80 en servicio, la exposición es 0,32.

Los tres siguientes dependen del cliente y de terceros. El riesgo R-11 se origina en que el cliente aporta 46 personas para nueve plataformas (Multitiendas Ancoa S.A., 2026, Caso 09, cap. 14, núm. 14.2, p. 32) y debe validar cada entregable en diez días hábiles. Cada validación tardía atrasa entre 3 y 6 semanas un entregable de la ruta crítica, por lo que rige la regla del hito contractual. Tiene probabilidad 0,40, impacto 0,80 y exposición 0,32.

El riesgo R-37 nace de que siete decisiones del cliente deben tomarse antes de cerrar el apartado 4.1. Una demora aplaza entre 3 y 6 semanas el cierre de la arquitectura frente a la línea base del mes 4. Por la regla del hito contractual el impacto es 0,80, y con probabilidad 0,40 la exposición es 0,32.

El riesgo R-20 se origina en que el cliente compra el hardware, contrata los enlaces y ejecuta las obras de 14 tiendas (Multitiendas Ancoa S.A., 2026, Caso 09, cap. 11, p. 27). Cada ola de despliegue depende de ello, y el atraso de una ola desplaza el despliegue gradual entre 3 y 6 semanas. Tiene probabilidad 0,60 e impacto 0,40 en plazo, con exposición 0,24.

Los tres siguientes se relacionan con el equipo y con el plan. El riesgo R-23 considera que, entre 85 profesionales y un núcleo de 16, la salida del jefe de proyecto, del arquitecto o del integrador a mitad de la Etapa 1 atrasaría entre 3 y 6 semanas la ruta crítica. Rige la regla del hito contractual, con impacto 0,80, probabilidad 0,40 y exposición 0,32.

El riesgo R-33 se origina en que, entre los meses 13 y 15, el mismo equipo sostiene la marcha blanca de la Etapa 1 y construye la Etapa 2, y la nivelación de recursos del Formulario T-15 cubre ese período solo en parte. La sobreasignación atrasa una de las dos ramas críticas, con hitos en los meses 16 y 21. Tiene probabilidad 0,40, impacto 0,80 y exposición 0,32.

El riesgo R-24 resulta de que la Etapa 1 entra en producción en abril de 2028 y el Día de la Madre y el evento anual llegan entre 4 y 10 semanas después. Los servicios de existencias y ventas enfrentan su primer peak con defectos latentes, con riesgo de repetir el episodio de junio de 2026, que dejó 2.840 cancelaciones (Multitiendas Ancoa S.A., 2026, Caso 09, cap. 7, núm. 7.1, p. 15). La probabilidad es 0,60, el impacto 0,40 y la exposición 0,24.

Los dos últimos se relacionan con aprobaciones y metas de operación. El riesgo R-16 se origina en que la autoridad tributaria no ha aprobado la modalidad fiscal de contingencia, según la Tabla 4.9 del capítulo 4. Sin ella, un corte de enlace en tienda incumple las restricciones 5 o 6 (Multitiendas Ancoa S.A., 2026, Caso 09, cap. 10, restricciones 5 y 6, p. 26). Tiene origen externo, probabilidad 0,40, impacto 0,80 y exposición 0,32.

El riesgo R-34 parte de que la meta de error de precio es menor que 3 \% y de que el papel no baja de 2,6 \% según el capítulo 2, con hasta 400.000 cambios de precio al día en campaña. El margen es estrecho y las oleadas diurnas de cambio de etiquetas dependen del personal de tienda. La probabilidad es 0,60, el impacto 0,40 en alcance y calidad y la exposición 0,24.

\subsection{Riesgos exigidos en zona media}

Ocho de los once riesgos exigidos quedan en zona media. En cada uno, la probabilidad es baja, el efecto se acota con un rodeo o con una salida parcial, o el impacto es moderado, y por eso no alcanzan la zona alta.

El riesgo R-03 tiene impacto 0,80, porque un cruce sin base quiebra la separación de los dos negocios que el Caso declara línea roja (Multitiendas Ancoa S.A., 2026, Caso 09, cap. 19, p. 46). La probabilidad es 0,20, porque el Servicio de control de cruces y el ADR-02 lo bloquean por diseño. La exposición es 0,16.

El riesgo R-04 parte de que nadie conoce hoy el mapa de las 14 interfaces (Multitiendas Ancoa S.A., 2026, Caso 09, cap. 7, núm. 7.4, p. 17). Su probabilidad es 0,60, pero el inventario del mes 4 acota el efecto a una degradación con rodeo, con impacto 0,20 y exposición 0,12.

El riesgo R-05 tiene probabilidad 0,85, porque la rotación anual del 62 \% (Multitiendas Ancoa S.A., 2026, Caso 09, cap. 7, núm. 7.4, p. 16) hace casi segura la pérdida de personal capacitado. El efecto es una degradación en caja con rodeos disponibles, con impacto 0,20 y exposición 0,17.

El riesgo R-06 considera que la comisión por canal de origen y la información precontractual alargan la venta con tarjeta desde el mes 16, antes de que el Servicio de comisiones llegue en el mes 21 (Multitiendas Ancoa S.A., 2026, Caso 09, cap. 8, entrevista al vendedor de piso, p. 21). Con probabilidad 0,40 e impacto 0,20 sobre áreas funcionales clave, queda en el límite inferior de la zona, con exposición 0,08.

El riesgo R-07 supone que en 56 meses es plausible que API Management Premium clásico, AKS o PostgreSQL lleguen a fin de soporte. La migración forzada costaría entre 1 y 2 \% de las horas, lo que da impacto 0,20. Con probabilidad 0,40, la exposición es 0,08.

El riesgo R-08 reconoce que salir de Event Hubs, API Management y AKS sería costoso, con impacto 0,40. Los clientes Kafka (ADR-06 y ADR-07) y la definición de la infraestructura como código dejan una salida parcial, de modo que la probabilidad es 0,20 y la exposición 0,08.

El riesgo R-09 parte de que el evento de 2026 concentró 104.000 pedidos y se proyectan 150.000 (Multitiendas Ancoa S.A., 2026, Caso 09, cap. 14, núm. 14.1, p. 31), con prueba a 225.000. La memoria de capacidad aún no está calculada y una saturación cancelaría pedidos como en junio. La probabilidad es 0,40, el impacto 0,40 y la exposición 0,16.

El riesgo R-10 tiene impacto 0,80, porque una fuga de datos de tarjeta con sanción es de efecto máximo, y la norma de medios de pago está definida solo para el canal digital (Multitiendas Ancoa S.A., 2026, Caso 09, cap. 6, p. 14). La probabilidad es 0,20, porque el compromiso exige vencer los controles de red de tienda y del emisor y existe prueba de penetración. La exposición es 0,16.

\subsection{Oportunidades}

El registro incluye ocho oportunidades, que se califican con la misma escala de probabilidad y con un impacto positivo. Las cuatro desarrolladas aquí son las de mayor efecto sobre los riesgos críticos y altos. Las otras cuatro (O-01, O-04, O-07 y O-08) se registran en el Formulario T-16. Se refieren al retiro anticipado de interfaces por archivo, a la entrada anticipada del Servicio de evidencia financiera, a la entrega temprana de documentación por los proveedores actuales y a la compra anticipada de hardware por el cliente. Ninguna oportunidad se cuenta en una reserva.

La oportunidad O-02 establece que, si la filial emisora fija los topes del cupo preaprobado antes del mes 6, entonces la prueba de corte de 24 h se adelanta y libera holgura para la migración de cartera. La oportunidad O-03 establece que, si la carga inicial de existencias usa el conteo cíclico y la detección de discrepancias de la segunda innovación, entonces la discrepancia del 12,4 \% baja antes del corte y se liberan horas de conteo en tienda. La oportunidad O-05 establece que, si la cartera sin casos difíciles migra completa antes del mes 22, entonces queda margen de observación antes del congelamiento de noviembre. La oportunidad O-06 establece que, si los datos del evento anual de 2026 y de ediciones previas calibran la memoria de capacidad antes del evento, entonces se dimensiona con medición y no con margen ciego, y luego se libera la capacidad no usada. Las probabilidades son 0,40 para O-02, O-03 y O-05 y 0,60 para O-06.

Los riesgos de adopción de las cinco innovaciones obligatorias no se repiten en este apartado y se tratan en el capítulo 13.

Con los riesgos identificados y calificados, el siguiente apartado presenta el análisis cuantitativo, que traduce esta exposición en semanas y horas de reserva para los riesgos críticos y altos.

\subsection{Análisis cuantitativo}

El análisis cuantitativo traduce la calificación del apartado anterior en semanas, horas y probabilidades, de modo que el plan pueda reservar capacidad con una base numérica. Se aplican dos técnicas porque responden preguntas distintas. La simulación del cronograma responde por el plazo del proyecto, esto es, cuánto atraso acumulado puede esperarse en los hitos contractuales. El análisis de modos de falla y efectos responde por la solución en operación, es decir, qué falla del sistema ya entregado merece más atención de diseño y de monitoreo. Todas las cifras son estimaciones del proponente que se recalculan cuando cambia un supuesto de la calificación.

\subsubsection{Valor esperado de los riesgos críticos y altos}

La primera técnica multiplica la probabilidad de cada uno de los 14 riesgos críticos y altos por las semanas y las horas que costaría su ocurrencia, antes de aplicar la respuesta. La Tabla 8.6 presenta el resultado. Los valores de semanas y horas son los del supuesto indicado en cada ficha del apartado 8.3, y los resultados se expresan solo en tiempo y esfuerzo.

Tabla 8.6: Valor esperado de los riesgos críticos y altos antes de la respuesta. Fuente: Elaboración propia.

\begin{longtable}{p{0.12\linewidth}p{0.12\linewidth}p{0.16\linewidth}p{0.16\linewidth}p{0.34\linewidth}}
\toprule
\textbf{Riesgo} & \textbf{Probabilidad} & \textbf{Semanas} & \textbf{Horas} & \textbf{Valor esperado} \\
\midrule
\endhead
\bottomrule
\endlastfoot
R-02 & 0,60 & 3 & 600 & 1,80 semanas y 360 horas \\
R-12 & 0,60 & 4 & 1.280 & 2,40 semanas y 768 horas \\
R-13 & 0,60 & 6 & 1.920 & 3,60 semanas y 1.152 horas \\
R-01 & 0,40 & 4 & 960 & 1,60 semanas y 384 horas \\
R-11 & 0,40 & 3 & 360 & 1,20 semanas y 144 horas \\
R-16 & 0,40 & 4 & 640 & 1,60 semanas y 256 horas \\
R-18 & 0,60 & 5 & 1.200 & 3,00 semanas y 720 horas \\
R-20 & 0,60 & 4 & 640 & 2,40 semanas y 384 horas \\
R-23 & 0,40 & 4 & 480 & 1,60 semanas y 192 horas \\
R-24 & 0,60 & 3 & 960 & 1,80 semanas y 576 horas \\
R-25 & 0,40 & 10 & 2.400 & 4,00 semanas y 960 horas \\
R-33 & 0,40 & 3 & 600 & 1,20 semanas y 240 horas \\
R-34 & 0,60 & 3 & 360 & 1,80 semanas y 216 horas \\
R-37 & 0,40 & 3 & 360 & 1,20 semanas y 144 horas \\
\textbf{Total} & & \textbf{59} & \textbf{12.760} & \textbf{29,20 semanas y 6.496 horas} \\
\end{longtable}

El valor esperado suma 29,2 semanas y 6.496 horas. Tres riesgos concentran más de un tercio de las semanas, a saber R-25, R-13 y R-18, lo que confirma que la migración de la cartera de crédito y la carga inicial de existencias son los frentes que más tiempo pueden consumir. Como las respuestas del apartado 8.3 se calculan sobre estos mismos riesgos, el total sirve de referencia para medir cuánto reducen la exposición.

\subsubsection{Análisis de modos de falla y efectos}

La segunda técnica examina cinco modos de falla de la solución en operación, elegidos porque combinan consecuencias graves con controles ya definidos en el diseño. Cada modo se puntúa de 1 a 10 en severidad (S), ocurrencia (O) y detección (D), donde 10 en detección significa que la falla pasa inadvertida. El número de prioridad de riesgo (NPR) es el producto de los tres puntajes. La severidad sigue al impacto de la calificación y la ocurrencia a su probabilidad. La Tabla 8.7 muestra el resultado.

Tabla 8.7: Análisis de modos de falla y efectos de la solución en operación. Fuente: Elaboración propia.

\begin{longtable}{p{0.30\linewidth}p{0.36\linewidth}p{0.18\linewidth}p{0.10\linewidth}}
\toprule
\textbf{Modo de falla} & \textbf{Efecto} & \textbf{S / O / D} & \textbf{NPR} \\
\midrule
\endhead
\bottomrule
\endlastfoot
Emisión tributaria en contingencia (R-16) & La tienda no emite documentos durante un corte. & 9 / 4 / 4 & 144 \\
Sincronización del diario local (R-15) & Cobro doble o descuadre con ERP y DTE. & 6 / 4 / 5 & 120 \\
Saturación en el evento anual (R-09) & Pedidos cancelados en los tres días de peak. & 7 / 4 / 4 & 112 \\
Cruce indebido de datos (R-03) & Se quiebra la separación entre retail y emisor. & 10 / 2 / 3 & 60 \\
Caída del centro de datos DC-01 (R-31) & Se detienen la emisión tributaria y el crédito. & 9 / 2 / 2 & 36 \\
\end{longtable}

La emisión en contingencia encabeza la lista porque combina una severidad alta, una ocurrencia moderada mientras la autoridad tributaria no apruebe la modalidad y una detección tardía. El cruce indebido de datos tiene la mayor severidad, pero su ocurrencia y su detección son bajas, porque el servicio de control de cruces lo bloquea por diseño y registra cada cruce; por eso queda bajo los demás en el NPR sin dejar de tener tolerancia cero. La caída de DC-01 se detecta de inmediato y es poco frecuente, aunque su efecto es grave y depende del recinto del cliente. En consecuencia, las pruebas de la sincronización, de la contingencia fiscal y de la capacidad del evento anual reciben la mayor atención de verificación antes del paso a producción.

\subsubsection{Simulación del cronograma}

La tercera lectura simula 10.000 veces el efecto conjunto de los mismos 14 riesgos sobre los hitos de los meses 16, 21 y 22. Cada riesgo se asigna a un eslabón de la ruta crítica del apartado 7.3.3, ocurre con su probabilidad y, si ocurre, aporta un atraso con distribución triangular entre la mitad y una vez y media las semanas estimadas. Los eventos se tratan como independientes. Dentro de un mismo eslabón los atrasos se solapan, porque corren en paralelo sobre la misma ventana, y se toma el mayor; entre eslabones distintos de la misma cadena se suman. El hito del mes 22 hereda la cadena de la Etapa 2 y agrega la migración de la cartera. Se corre el escenario sin respuestas, con las probabilidades originales, y el escenario con respuestas, con las probabilidades residuales del apartado 8.3. La Tabla 8.8 reúne los percentiles en semanas para el caso base solapado.

Tabla 8.8: Atraso acumulado por hito, sin y con respuestas. Fuente: Elaboración propia.

\begin{longtable}{p{0.34\linewidth}p{0.18\linewidth}p{0.12\linewidth}p{0.12\linewidth}p{0.12\linewidth}}
\toprule
\textbf{Hito} & \textbf{Escenario} & \textbf{P50} & \textbf{P80} & \textbf{P90} \\
\midrule
\endhead
\bottomrule
\endlastfoot
Mes 16, paso a producción de la Etapa 1 & Sin respuestas & 13,8 & 17,5 & 19,4 \\
Mes 16, paso a producción de la Etapa 1 & Con respuestas & 8,7 & 12,5 & 14,6 \\
Mes 21, paso a producción de la Etapa 2 & Sin respuestas & 0,0 & 3,0 & 3,4 \\
Mes 21, paso a producción de la Etapa 2 & Con respuestas & 0,0 & 1,7 & 3,0 \\
Mes 22, retiro de la plataforma de 2011 & Sin respuestas & 7,8 & 11,3 & 13,0 \\
Mes 22, retiro de la plataforma de 2011 & Con respuestas & 4,8 & 8,5 & 10,6 \\
\end{longtable}

Los hitos son contractuales y no se mueven, de modo que el atraso simulado no es una fecha nueva sino la cantidad de trabajo que el plan debe absorber. Por eso el P80 con respuestas se adopta como reserva de contingencia en semanas, que el apartado 8.3 reparte mediante fechas internas más tempranas. Esa reserva es de 12,5 semanas para el mes 16, 1,7 para el mes 21 y 8,5 para el mes 22. El hito del mes 16 concentra la mayor exposición porque seis eslabones de su cadena están expuestos, y las respuestas reducen su P80 en cerca de 5 semanas. Como cota superior prudente se informa además la suma en serie de todos los atrasos, que eleva el P80 con respuestas a 15,3 semanas en el mes 16 y a 10,0 en el mes 22, y que se interpreta como límite y no como valor esperado de planificación.

\section{8.3 Plan de Acción a Riesgos}\label{plan-de-acciuxf3n-a-riesgos}

Este apartado define la respuesta a cada riesgo del registro. Los riesgos críticos y altos reciben acciones concretas que viven en la estructura de desglose del trabajo, con responsable, disparador y plan de contingencia, y los riesgos medios y bajos se aceptan de forma activa o pasiva según la regla de zona del apartado 8.1. El detalle de los 37 riesgos está en el anexo, en las Tablas A.1 a A.3, y el de las oportunidades en la Tabla A.4. Los totales de semanas y horas que cierran el apartado alimentan las reservas de contingencia y de gestión.

\subsection{Estrategias de respuesta}

La estrategia de cada riesgo se eligió según su zona y según la causa que lo origina. Cuando la causa depende del diseño o del calendario del proponente, se prefirió evitar, porque una decisión temprana elimina el evento. Es el caso de R-12, que se evita al certificar en Preproducción y programar el despliegue fuera de los congelamientos, y de R-13, R-25 y R-37, que se evitan al fijar secuencias y fechas antes de que el atraso ocurra. Cuando el evento no puede eliminarse pero sí su probabilidad o su efecto, se mitigó, como en R-02, R-01, R-11, R-18, R-23, R-24, R-33 y R-34. R-20 se transfiere al cliente mediante obligaciones contractuales con fechas por ola, porque el hardware, los enlaces y las obras son de su responsabilidad. R-16 se escala, porque la aprobación de la modalidad fiscal de contingencia corresponde a la autoridad tributaria y el proponente solo puede gestionarla. Se acepta únicamente lo que la regla de zona permite aceptar, esto es, riesgos medios y bajos, de los cuales R-19, R-29 y R-38 se aceptan con monitoreo y los demás se mitigan. Para las oportunidades se usan las estrategias simétricas, que son explotar, mejorar y compartir. La Tabla 8.9 resume las respuestas de los 14 riesgos críticos y altos.

Tabla 8.9: Respuesta a los riesgos críticos y altos. Fuente: Elaboración propia.

\begin{longtable}{p{0.069\linewidth}p{0.109\linewidth}p{0.406\linewidth}p{0.178\linewidth}p{0.168\linewidth}}
\toprule
\textbf{ID} & \textbf{Estrategia} & \textbf{Acción} & \textbf{Cuenta de la EDT} & \textbf{Responsable} \\
\midrule
\endhead
\bottomrule
\endlastfoot
R-02 & Mitigar & Declarar una reserva en semanas y no desplegar servicios nuevos en las ocho semanas previas al evento. & 1.1.6, 1.9.6 & Jefe de Proyecto \\
R-12 & Evitar & Certificar en Preproducción y desplegar en la ventana libre de enero de 2028. & 1.12.1, 1.1.6 & Jefe de Proyecto \\
R-13 & Evitar & Fijar la secuencia de cortes por tramo y cerrar el último antes del 15 de octubre de 2028. & 1.7.10, 1.7.11, 1.7.14 & Líder de Datos y Migración \\
R-01 & Mitigar & Cortar por tramo con conciliación diaria y criterio de detención acordado con la filial. & 1.7.1, 1.7.7, 1.7.14, 1.5.11.5 & Líder de Datos y Migración \\
R-11 & Mitigar & Pactar actas con plazo de diez días hábiles y escalar la validación vencida. & 1.1.7, 1.1.10 & Jefe de Proyecto \\
R-16 & Escalar & Gestionar la aprobación fiscal y escalarla al Comité Ejecutivo si no avanza. & 1.6.12 & Jefe de Proyecto \\
R-18 & Mitigar & Cargar existencias con conteo cíclico y conciliar las 24 instalaciones antes del corte. & 1.7.3, 1.7.4, 1.5.3.4 & Líder de Datos y Migración \\
R-20 & Transferir & Fijar por contrato las responsabilidades del cliente y exigir certificación de sitio por ola. & 1.3.8, 1.4.8, 1.11.3 & Líder de Implantación \\
R-23 & Mitigar & Nombrar suplentes, documentar el conocimiento y adelantar la contratación. & 1.1.1 & Jefe de Proyecto \\
R-24 & Mitigar & Ampliar el soporte reforzado y cerrar los defectos graves antes de abril de 2028. & 1.9.6, 1.12.9 & Líder de Operación \\
R-25 & Evitar & Cortar el último tramo antes del 15 de octubre y acordar un plan de retorno. & 1.7.10, 1.7.11 & Líder de Datos y Migración \\
R-33 & Mitigar & Nivelar los recursos de los meses 13 a 15 y reforzar el equipo. & Nivelación de recursos (T-15) & Jefe de Proyecto \\
R-34 & Mitigar & Dimensionar las oleadas de cambio de etiquetas y evaluar etiquetas electrónicas. & 1.5.1.2, 1.2.7 & Líder de Operación \\
R-37 & Evitar & Calendarizar las siete decisiones del cliente antes de la línea base y escalar las vencidas. & 1.2.6, 1.1.4, 1.1.3 & Jefe de Proyecto \\
\end{longtable}

Ocho de las catorce respuestas son mitigaciones, cuatro evitan el riesgo, una lo transfiere y una lo escala. Siete tienen como responsable al Jefe de Proyecto y cuatro al Líder de Datos y Migración, de modo que la ejecución del plan se concentra en los roles que ya sostienen la ruta crítica.

\subsection{Respuesta a los riesgos críticos}

Para R-02, la respuesta consiste en declarar y cuantificar una reserva en semanas para la fecha del evento anual, tomando fechas internas más tempranas dentro de las ventanas existentes, y en no desplegar servicios nuevos dentro de las ocho semanas previas al evento. El disparador es recibir la fecha con menos de ocho semanas de aviso o tener un servicio a medio desplegar en esa fecha. Si la respuesta falla, el Jefe de Proyecto reordena las entregas con cinco personas durante tres semanas, lo que equivale a 600 horas. La respuesta baja la probabilidad de 0,60 a 0,40, pero el impacto se mantiene en 0,80 por la regla del hito contractual, de modo que la exposición residual es 0,32 y el riesgo queda en zona alta, bajo seguimiento del Comité Ejecutivo.

Para R-12, la respuesta evita el choque con los congelamientos. La certificación de la Etapa 1 se hace en Preproducción en el mes 12, que no es un sistema en producción del cliente. El despliegue que inicia la marcha blanca se ejecuta en la ventana libre del 7 al 24 de enero de 2028, solo con lo certificado, y las correcciones se liberan en la ventana libre que va desde la segunda semana de marzo hasta fines de ese mes, unas tres semanas y media antes del paso a producción de abril. Durante los congelamientos solo se admiten correcciones de emergencia autorizadas por el Comité Ejecutivo. El disparador es tener defectos de severidad alta abiertos al 20 de enero de 2028 o el despliegue sin terminar al 24 de enero. Si la respuesta falla y se pierde la ventana de corrección, el Jefe de Proyecto solicita al Comité Ejecutivo las correcciones de emergencia y reasigna ocho personas durante cuatro semanas, equivalentes a 1.280 horas. La probabilidad residual es 0,20, la exposición 0,16 y la zona media. Las fechas dependen de SUP-26 y se recalculan si cambia.

Para R-13, la respuesta fija desde el mes 16 la secuencia de cortes por tramo, con el último corte antes del 15 de octubre de 2028, reserva el tramo complejo para la marcha blanca de la Etapa 2 y acuerda con la filial emisora un criterio de detención por discrepancia. El disparador es que algún tramo no tenga corte al 15 de septiembre de 2028. Si la respuesta falla, se rehace el corte del tramo y su conciliación con ocho personas durante seis semanas, que suman 1.920 horas. La ventana de migración deja una holgura casi nula, por lo que la probabilidad solo baja de 0,60 a 0,40 y la exposición residual es 0,32. R-13 conserva así exposición alta tras la respuesta y se escala al Comité Ejecutivo, que decide ante cualquier rebase del plazo anterior a 2029.

\subsection{Respuesta a los riesgos altos}

Los once riesgos altos se agrupan por tema. En la cartera y el retiro de la plataforma de 2011, R-01 se mitiga cortando por tramo con conciliación diaria y reservando el tramo difícil para la marcha blanca de la Etapa 2, y su disparador es una divergencia sin conciliar dos días seguidos en un tramo. R-25 se evita cortando el último tramo antes del 15 de octubre de 2028 y acordando con el Comité Ejecutivo un plan de retorno para el congelamiento, y su disparador es cualquier tramo sin corte en esa fecha. R-18 se mitiga con carga por conteo cíclico y conciliación de las 24 instalaciones, y su disparador es una discrepancia de la carga de prueba sobre el umbral acordado.

Para el primer peak tras la producción, R-24 amplía el soporte reforzado entre el paso a producción y el Día de la Madre y el evento anual, y su disparador es tener defectos de severidad alta abiertos a cuatro semanas del Día de la Madre. R-34 dimensiona las oleadas de cambio de etiquetas para sostener la meta de menos de 3 \% y evalúa las etiquetas electrónicas, y su disparador es un error de precio sobre 3 \% en el piloto.

Las dependencias del cliente se tratan de dos maneras. R-11 y R-37 se resuelven con actas de plazo fijo y con un calendario de las siete decisiones antes de la línea base del mes 4, y sus disparadores son un acta sin validar al décimo día hábil y decisiones pendientes al mes 3. R-20 se transfiere con obligaciones contractuales por ola y su disparador es una certificación de sitio no emitida ocho semanas antes de la ola. R-16 se escala y su disparador es la modalidad sin aprobar al mes 9, antes de la certificación del mes 12.

En el equipo del proponente, R-23 se mitiga con suplentes y traspaso documentado, con disparador en una vacante de rol clave sin cubrir por más de 30 días. R-33 se mitiga nivelando los recursos de los meses 13 a 15, con disparador en una carga sobre 100 \% de un rol clave. La Tabla 8.10 presenta el disparador, el plan de contingencia y la exposición antes y después de la respuesta.

Tabla 8.10: Disparadores y contingencia de los riesgos altos. Fuente: Elaboración propia.

\begin{longtable}{p{0.066\linewidth}p{0.275\linewidth}p{0.323\linewidth}p{0.133\linewidth}p{0.133\linewidth}}
\toprule
\textbf{ID} & \textbf{Disparador} & \textbf{Plan de contingencia} & \textbf{Exposición antes} & \textbf{Exposición residual} \\
\midrule
\endhead
\bottomrule
\endlastfoot
R-01 & Divergencia sin conciliar dos días seguidos. & Repetir el corte del tramo con seis personas. & 0,32 & 0,16 \\
R-11 & Acta sin validar al décimo día hábil. & Escalar a la Contraparte Técnica. & 0,32 & 0,16 \\
R-16 & Modalidad sin aprobar al mes 9. & Rediseñar la emisión sin enlace con cuatro personas. & 0,32 & 0,16 \\
R-18 & Discrepancia de la carga sobre el umbral. & Repetir la carga y la conciliación con seis personas. & 0,24 & 0,08 \\
R-20 & Certificación de sitio no emitida a ocho semanas. & Reprogramar la ola con cuatro personas. & 0,24 & 0,08 \\
R-23 & Vacante clave sin cubrir por 30 días. & Reemplazar y traspasar el rol con tres personas. & 0,32 & 0,16 \\
R-24 & Defectos graves abiertos a cuatro semanas. & Corregir en producción con ocho personas. & 0,24 & 0,08 \\
R-25 & Tramo sin corte al 15 de octubre. & Activar el plan de retorno acordado. & 0,32 & 0,16 \\
R-33 & Carga sobre 100 \% en los meses 13 a 15. & Reasignar cinco personas entre las ramas. & 0,32 & 0,16 \\
R-34 & Error de precio sobre 3 \% en el piloto. & Rediseñar el proceso de etiquetas con tres personas. & 0,24 & 0,08 \\
R-37 & Decisiones pendientes al mes 3. & Reasignar a tres personas de arquitectura mientras se decide. & 0,32 & 0,16 \\
\end{longtable}

La respuesta reduce la exposición de los once riesgos altos a la mitad o a un tercio y los lleva a zona media. Cada plan de contingencia dimensiona el esfuerzo de recuperación con el mismo supuesto que sostiene la estimación de horas de la Tabla 8.6, de modo que la contingencia y la reserva se calculan con una base común.

\subsection{Costo-beneficio de las respuestas}

El beneficio de cada respuesta se mide como la reducción del valor esperado en horas, esto es, la probabilidad por las horas que costaría el evento, antes y después de la acción. El valor esperado residual usa la probabilidad residual de cada riesgo con las mismas horas. Los antecedentes del plan no estiman el esfuerzo de ejecutar cada acción, por lo que no se calcula una relación entre beneficio y esfuerzo y el beneficio se expresa solo como reducción. La Tabla 8.11 presenta el resultado.

Tabla 8.11: Reducción del valor esperado por la respuesta. Fuente: Elaboración propia.

\begin{longtable}{p{0.099\linewidth}p{0.218\linewidth}p{0.218\linewidth}p{0.218\linewidth}p{0.178\linewidth}}
\toprule
\textbf{ID} & \textbf{Antes (horas)} & \textbf{Residual (horas)} & \textbf{Reducción (horas)} & \textbf{Reducción (\%)} \\
\midrule
\endhead
\bottomrule
\endlastfoot
R-02 & 360 & 240 & 120 & 33 \\
R-12 & 768 & 256 & 512 & 67 \\
R-13 & 1.152 & 768 & 384 & 33 \\
R-01 & 384 & 192 & 192 & 50 \\
R-11 & 144 & 72 & 72 & 50 \\
R-16 & 256 & 128 & 128 & 50 \\
R-18 & 720 & 480 & 240 & 33 \\
R-20 & 384 & 256 & 128 & 33 \\
R-23 & 192 & 96 & 96 & 50 \\
R-24 & 576 & 384 & 192 & 33 \\
R-25 & 960 & 480 & 480 & 50 \\
R-33 & 240 & 120 & 120 & 50 \\
R-34 & 216 & 144 & 72 & 33 \\
R-37 & 144 & 72 & 72 & 50 \\
\textbf{Total} & \textbf{6.496} & \textbf{3.688} & \textbf{2.808} & \textbf{43} \\
\end{longtable}

Las respuestas reducen el valor esperado de 6.496 a 3.688 horas, es decir, en 2.808 horas o 43 \% (2.808 dividido por 6.496). Cuatro respuestas aportan 1.616 de esas horas, equivalentes a 58 \% de la reducción. Son la de R-12, con 512 horas, la de R-25, con 480, la de R-13, con 384, y la de R-18, con 240. R-12 rinde más porque cambiar el calendario de despliegue baja su probabilidad de 0,60 a 0,20, y R-25 y R-13 rinden por las horas elevadas que cada evento costaría. Las respuestas de R-02, R-13, R-18, R-20, R-24 y R-34 reducen la probabilidad solo de 0,60 a 0,40, por lo que su reducción es de un tercio y el riesgo conserva un peso relevante. La atención del Comité de Proyecto debe concentrarse en esos seis.

\subsection{Cuentas de control propuestas al plan de trabajo}

Cuatro respuestas no tienen todavía una cuenta de control donde ejecutarse, y por eso se proponen al plan de trabajo del capítulo 7. La cuenta de control propuesta «Acuerdo con la filial emisora sobre la regla de comisión del vendedor» cubre R-06. La cuenta propuesta «Acuerdos de cooperación y salida con los proveedores de las plataformas que se reemplazan» cubre R-27. La cuenta propuesta «Medida transitoria de retención de grabaciones de repactación» cubre R-28, y la cuenta propuesta «Plan de incorporación de transportistas y vendedores de marketplace a los contratos de estado» cubre R-36. Cada una tiene su ventana definida antes del hito en que el riesgo actúa. Ninguna debería integrar la ruta crítica, salvo que el acuerdo con los proveedores demore la línea base del mes 4, caso en que se vigila como parte de R-27.

\subsection{Reserva de contingencia}

La reserva de contingencia cubre los riesgos identificados y se dimensiona con el percentil 80 de la simulación con respuestas, en el caso base solapado de la Tabla 8.8. En plazo, la reserva es de 12,5 semanas antes del hito del mes 16, 1,7 semanas antes del hito del mes 21 y 8,5 semanas antes del hito del mes 22. Como cota superior prudente, la suma en serie eleva esas cifras a 15,3, 1,7 y 10,0 semanas, y se informa para que el Comité de Proyecto conozca el límite.

Los hitos son contractuales y los eslabones principales de la ruta crítica tienen holgura cero, por lo que la reserva no alarga el plazo. El plan de trabajo ya fija fechas internas anteriores a las anclas, como la versión candidata de la Etapa 1 verificada antes del 1 de noviembre, al terminar el mes 10, y reconoce un margen de tres meses en el plan de marcha blanca y en la capacitación de administradores de la Etapa 1, y de un mes en los de la Etapa 2. Se refleja en el cronograma como fechas internas de término más tempranas, distribuidas por eslabón de la ruta crítica, a saber el levantamiento, la arquitectura, la plataforma base, el desarrollo, la certificación y la marcha blanca. Cada eslabón recibe una parte proporcional al atraso esperado de los riesgos que se le asignaron. Además, la reserva se materializa en ventanas de corrección ya identificadas fuera de los congelamientos, como la de enero de 2028 para el despliegue y la de marzo para las correcciones. Las 12,5 semanas del mes 16 equivalen a 19 \% de los 15 meses de la Etapa 1, que son 65 semanas (15 meses por 52 semanas, dividido por 12 meses; 12,5 dividido por 65). Es una exigencia alta sobre el plan, y por eso se vigila en cada revisión del registro.

En esfuerzo, la reserva es la suma del valor esperado residual de los 14 riesgos, esto es, 3.688 horas. La autoriza el Jefe de Proyecto, que informa cada uso al Comité de Proyecto en su siguiente sesión.

\subsection{Reserva de gestión}

La reserva de gestión cubre los riesgos no identificados, es decir, los que el registro no contiene. Su liberación corresponde únicamente al Comité Ejecutivo, como lo establece el umbral de escalamiento del apartado 8.1. Su valor, y el reflejo de ambas reservas en el flujo de caja, se informan únicamente en la oferta económica, por lo que este capítulo no contiene cifras monetarias.

\subsection{Seguimiento y oportunidades}

Los disparadores se vigilan en el ciclo de revisión del apartado 8.1. El registro se revisa cada quince días en el Comité de Proyecto y se actualiza cuando un disparador se activa, cuando una respuesta se completa o cuando cambia un supuesto como SUP-26 o SUP-28. Cada uso de la reserva de contingencia y cada riesgo nuevo quedan anotados con su fecha, su responsable y su efecto en la exposición.

Las oportunidades se gestionan con las mismas instancias, pero no se suman a las reservas. La filial emisora puede fijar los topes del cupo preaprobado antes del mes 6, lo que adelanta la prueba de corte de 24 horas (O-02, mejorar, a cargo del Jefe de Proyecto). El conteo cíclico de la innovación 2 puede bajar la discrepancia de 12,4 \% antes del corte (O-03, explotar, a cargo del Líder de Datos y Migración). La cartera sin casos difíciles puede migrar completa antes del mes 22, lo que deja margen de observación antes del congelamiento de noviembre (O-05, mejorar, con el mismo responsable). Los datos del evento anual de 2026 permiten calibrar la memoria de capacidad con medición (O-06, mejorar, a cargo del Líder de Rendimiento y Calidad). En conjunto, las respuestas reducen en 43 \% el valor esperado de los riesgos críticos y altos, las reservas se dimensionan con la simulación y el seguimiento mantiene vigentes los disparadores hasta el cierre del contrato.

\section*{Anexo A · Registro de riesgos (Formulario T-16)}
\addcontentsline{toc}{section}{Anexo A · Registro de riesgos (Formulario T-16)}

Este anexo reúne los 37 riesgos y las 8 oportunidades del registro de Only Simple Solutions para el proyecto de Multitiendas Ancoa S.A. Los campos siguen la práctica de un registro de riesgos, con identificación, análisis y respuesta, y están sujetos a la estructura definitiva del Formulario T-16. La Tabla A.1 identifica cada riesgo, la Tabla A.2 presenta su análisis, la Tabla A.3 define su respuesta y la Tabla A.4 registra las oportunidades.

La Tabla A.1 identifica los 37 riesgos con su enunciado resumido, su rama de la estructura de desglose y la mirada desde la que se observan.

Tabla A.1: Identificación de los riesgos. Fuente: Elaboración propia.

\begin{longtable}{p{0.080\linewidth}p{0.540\linewidth}p{0.160\linewidth}p{0.160\linewidth}}
\toprule
\textbf{ID} & \textbf{Riesgo} & \textbf{Rama} & \textbf{Mirada} \\
\midrule
\endhead
\bottomrule
\endlastfoot
R-01 & Los saldos de la cartera migrados por tramos divergen de la plataforma de 2011 y se incumple la restricción 8. & Técnica & Implantación \\
R-02 & El evento anual cae con un servicio a medio desplegar o dentro de una marcha blanca. & Proyecto & Implantación \\
R-03 & Un dato del negocio financiero llega al retail sin base ni registro, y se quiebra la separación de los dos negocios. & Seguridad & Solución \\
R-04 & Aparece en producción una integración no inventariada entre las 14 interfaces y falla una operación. & Técnica & Desarrollo, Implantación \\
R-05 & La rotación del 62 \% hace que el personal capacitado se vaya antes o durante la marcha blanca. & Organizacional & Implantación \\
R-06 & La comisión por canal y la información precontractual alargan la venta con tarjeta y el vendedor deja de ofrecerla. & Organizacional & Solución, Implantación \\
R-07 & Un servicio administrado llega a fin de soporte durante los 56 meses y exige una migración forzada. & Técnica & Solución \\
R-08 & El uso de servicios propios de Azure ata a Ancoa a un proveedor y encarece cambiarse. & Técnica & Solución \\
R-09 & La reserva de existencia o la entrada API se saturan en el peak del evento anual y se cancelan pedidos. & Técnica & Solución \\
R-10 & Un compromiso de seguridad en tienda o en el emisor produce una fuga de datos de tarjeta y una sanción. & Seguridad & Solución \\
R-11 & La Contraparte Técnica o la filial no validan dentro de diez días hábiles y se atrasan entregables de la ruta crítica. & Organizacional & Desarrollo \\
R-12 & La marcha blanca de los meses 13 a 15 cae en los congelamientos y arrastra el paso a producción del mes 16. & Proyecto & Implantación \\
R-13 & Un tramo de la cartera no alcanza a migrar antes del retiro de la plataforma de 2011 en el mes 22. & Proyecto & Desarrollo, Implantación \\
R-14 & Sin topes del cupo o con la prueba de corte de 24 h fallida, cae la venta con crédito en contingencia. & Proyecto & Desarrollo \\
R-15 & La sincronización del diario local tras un corte duplica o pierde ventas en las tiendas sin enlace de respaldo. & Técnica & Solución \\
R-16 & La autoridad tributaria no aprueba la modalidad fiscal de contingencia y la tienda no puede emitir durante un corte. & Proyecto & Solución \\
R-17 & La CMF o la asesoría jurídica objetan la recuperación del emisor en Brazil South y obligan a rediseñarla. & Proyecto & Solución \\
R-18 & La carga inicial de existencias hereda los errores del registro y el disponible publicado no es confiable en el mes 16. & Técnica & Implantación \\
R-19 & La reglamentación de la Ley 21.719 precisa controles no previstos y exige retrabajo en cruces y evidencia. & Proyecto & Solución \\
R-20 & El hardware, los enlaces o las obras que provee el cliente no están listos para la ola de una tienda. & Proyecto & Implantación \\
R-21 & El cliente rechaza un supuesto de los 25 y el alcance cambia después de la línea base del mes 4. & Proyecto & Desarrollo \\
R-22 & Los repositores externos y las terminales compartidas impiden asegurar acceso y los conteos quedan incompletos. & Operación & Solución, Implantación \\
R-23 & Sale un rol clave del equipo o la contratación no llega a tiempo, y se atrasa la ruta crítica. & Organizacional & Desarrollo \\
R-24 & Los servicios de existencias y ventas enfrentan su primer peak con defectos latentes, como en junio de 2026. & Operación & Implantación \\
R-25 & Una falla posterior al retiro de la plataforma de 2011 queda sin corrección hasta enero de 2029. & Proyecto & Implantación \\
R-27 & Los seis proveedores actuales no documentan interfaces ni sostienen la convivencia, y se atrasan integración y migración. & Proyecto & Desarrollo \\
R-28 & Las grabaciones se borran a los 90 días hasta el mes 16 y crece el pasivo de las 1.240 repactaciones sin evidencia. & Proyecto & Desarrollo \\
R-29 & El procedimiento abierto por el SERNAC impone obligaciones nuevas y cambia el alcance después de la línea base. & Proyecto & Desarrollo \\
R-30 & Un servicio Premium no puede crearse en Chile Central y obliga a rediseñar o atrasa la plataforma base. & Técnica & Desarrollo \\
R-31 & Un evento local o una obra no lista en el centro de datos del cliente detiene la emisión tributaria o el crédito. & Operación & Solución \\
R-32 & El inventario cambia la ubicación de las plataformas y los plazos de los enlaces, y se atrasa la plataforma base. & Proyecto & Desarrollo \\
R-33 & El mismo equipo sostiene la marcha blanca de la Etapa 1 y construye la Etapa 2 en los meses 13 a 15. & Proyecto & Desarrollo \\
R-34 & El personal de tienda no alcanza a cambiar las etiquetas de papel y no se cumple la meta de error de precio. & Operación & Solución \\
R-35 & Quedan accesos de ex trabajadores o temporales sin revocar y se quiebra la segregación. & Seguridad & Solución \\
R-36 & Transportistas y vendedores de marketplace no adoptan a tiempo los contratos de estado y la meta del 97 \% no se mide. & Proyecto & Implantación \\
R-37 & Las siete decisiones del cliente se demoran y la arquitectura no cierra con la línea base del mes 4. & Organizacional & Desarrollo \\
R-38 & Una falla del piloto o de una ola en el clúster de Santiago concentra el impacto en la venta. & Operación & Implantación \\
\end{longtable}

La tabla muestra que cada rama de la estructura contiene al menos un riesgo y que cada riesgo se observa desde una o más miradas. La Tabla A.2 presenta la calificación de cada uno.

Tabla A.2: Análisis cualitativo de los riesgos. Fuente: Elaboración propia.

\begin{longtable}{p{0.12\linewidth}p{0.18\linewidth}p{0.18\linewidth}p{0.20\linewidth}p{0.20\linewidth}}
\toprule
\textbf{ID} & \textbf{Probabilidad} & \textbf{Impacto} & \textbf{Exposición} & \textbf{Zona} \\
\midrule
\endhead
\bottomrule
\endlastfoot
R-01 & 0,40 & 0,80 & 0,32 & Alta \\
R-02 & 0,60 & 0,80 & 0,48 & Crítica \\
R-03 & 0,20 & 0,80 & 0,16 & Media \\
R-04 & 0,60 & 0,20 & 0,12 & Media \\
R-05 & 0,85 & 0,20 & 0,17 & Media \\
R-06 & 0,40 & 0,20 & 0,08 & Media \\
R-07 & 0,40 & 0,20 & 0,08 & Media \\
R-08 & 0,20 & 0,40 & 0,08 & Media \\
R-09 & 0,40 & 0,40 & 0,16 & Media \\
R-10 & 0,20 & 0,80 & 0,16 & Media \\
R-11 & 0,40 & 0,80 & 0,32 & Alta \\
R-12 & 0,60 & 0,80 & 0,48 & Crítica \\
R-13 & 0,60 & 0,80 & 0,48 & Crítica \\
R-14 & 0,40 & 0,20 & 0,08 & Media \\
R-15 & 0,40 & 0,40 & 0,16 & Media \\
R-16 & 0,40 & 0,80 & 0,32 & Alta \\
R-17 & 0,40 & 0,40 & 0,16 & Media \\
R-18 & 0,60 & 0,40 & 0,24 & Alta \\
R-19 & 0,40 & 0,20 & 0,08 & Media \\
R-20 & 0,60 & 0,40 & 0,24 & Alta \\
R-21 & 0,60 & 0,20 & 0,12 & Media \\
R-22 & 0,60 & 0,20 & 0,12 & Media \\
R-23 & 0,40 & 0,80 & 0,32 & Alta \\
R-24 & 0,60 & 0,40 & 0,24 & Alta \\
R-25 & 0,40 & 0,80 & 0,32 & Alta \\
R-27 & 0,40 & 0,40 & 0,16 & Media \\
R-28 & 0,85 & 0,20 & 0,17 & Media \\
R-29 & 0,20 & 0,20 & 0,04 & Baja \\
R-30 & 0,20 & 0,40 & 0,08 & Media \\
R-31 & 0,20 & 0,80 & 0,16 & Media \\
R-32 & 0,40 & 0,20 & 0,08 & Media \\
R-33 & 0,40 & 0,80 & 0,32 & Alta \\
R-34 & 0,60 & 0,40 & 0,24 & Alta \\
R-35 & 0,20 & 0,80 & 0,16 & Media \\
R-36 & 0,40 & 0,20 & 0,08 & Media \\
R-37 & 0,40 & 0,80 & 0,32 & Alta \\
R-38 & 0,20 & 0,20 & 0,04 & Baja \\
\end{longtable}

La exposición es el producto de la probabilidad por el impacto y la zona se determina según el apartado 8.1. La Tabla A.3 define la respuesta a cada riesgo.

Tabla A.3: Respuesta a los riesgos. Fuente: Elaboración propia.

\begin{longtable}{p{0.065\linewidth}p{0.102\linewidth}p{0.316\linewidth}p{0.158\linewidth}p{0.288\linewidth}}
\toprule
\textbf{ID} & \textbf{Estrategia} & \textbf{Acción} & \textbf{Responsable} & \textbf{Disparador} \\
\midrule
\endhead
\bottomrule
\endlastfoot
R-01 & Mitigar & Fijar cortes por tramo con conciliación diaria y criterio de detención acordado con la filial. & Líder de Datos y Migración & Divergencia sin conciliar dos días seguidos en un tramo. \\
R-02 & Mitigar & Reservar semanas para la fecha del evento y no desplegar servicios nuevos en las ocho semanas previas. & Jefe de Proyecto & Fecha del evento con menos de 8 semanas de aviso. \\
R-03 & Evitar & Mantener el Servicio de control de cruces como único camino de datos, con base legal por cruce. & Arquitecto de Soluciones & Un solo cruce sin base o sin registro. \\
R-04 & Mitigar & Cerrar el inventario antes de la línea base y certificar cada integración antes de retirar una plataforma. & Líder de Integración & Inventario no cerrado en el mes 4. \\
R-05 & Mitigar & Exigir la capacitación como condición de cierre de la marcha blanca y repetirla en temporada. & Líder de Gestión del Cambio & Personal capacitado bajo el umbral de cierre. \\
R-06 & Mitigar & Acordar con la filial la regla de comisión antes del mes 13 y simplificar la información precontractual. & Jefe de Proyecto & Acta de comisión sin firma en el mes 13. \\
R-07 & Mitigar & Llevar un calendario de ciclo de vida y planificar cada migración con un año de anticipación. & Arquitecto de Soluciones & Servicio a menos de 12 meses de su fin de soporte. \\
R-08 & Mitigar & Definir la infraestructura como código y un plan de reversibilidad con clientes Kafka. & Arquitecto de Soluciones & Plan de reversibilidad no aceptado en su fecha. \\
R-09 & Mitigar & Calcular la memoria de capacidad con datos de 2026 y probar a 225.000 pedidos. & Líder de Rendimiento y Calidad & Prueba a 225.000 pedidos sin cumplir antes del mes 16. \\
R-10 & Mitigar & Delimitar el ámbito de medios de pago, segmentar la red de tienda y ejecutar prueba de penetración. & Responsable de Seguridad & Hallazgo crítico en la prueba de penetración. \\
R-11 & Mitigar & Pactar actas con plazo de diez días hábiles y escalar la validación vencida. & Jefe de Proyecto & Acta sin validar al décimo día hábil. \\
R-12 & Evitar & Certificar en el mes 12, desplegar en la ventana libre de enero de 2028 y limitar el congelamiento a emergencias. & Jefe de Proyecto & Defectos altos abiertos al 20 de enero de 2028. \\
R-13 & Evitar & Fijar la secuencia de cortes por tramo con último corte antes del 15 de octubre de 2028 y escalar al Comité Ejecutivo. & Líder de Datos y Migración & Algún tramo sin corte al 15 de septiembre de 2028. \\
R-14 & Mitigar & Fijar con la filial los topes antes del mes 6 y mantener otro medio de pago como plan B. & Arquitecto de Soluciones & Topes sin fijar en el mes 6 o prueba de 24 h fallida. \\
R-15 & Mitigar & Dimensionar según el inventario de enlaces y probar la sincronización tras un corte de 24 h. & Arquitecto de Soluciones & Prueba de 24 h con una venta duplicada o perdida. \\
R-16 & Escalar & Gestionar la aprobación con la autoridad tributaria y escalar al Comité Ejecutivo si no avanza. & Jefe de Proyecto & Modalidad sin aprobar en el mes 9. \\
R-17 & Mitigar & Obtener el pronunciamiento jurídico y de la CMF y mantener una alternativa de recuperación en Chile. & Responsable de Seguridad & Sin pronunciamiento antes de la línea base del mes 4. \\
R-18 & Mitigar & Cargar con conteo cíclico y detección de discrepancias y conciliar las 24 instalaciones antes del corte. & Líder de Datos y Migración & Discrepancia de la carga de prueba sobre el umbral. \\
R-19 & Aceptar & Monitorear la reglamentación y absorber los ajustes en la matriz de cumplimiento. & Responsable de Seguridad & Publicación de una norma que cambie un control. \\
R-20 & Transferir & Fijar por contrato las responsabilidades del cliente con fechas por ola y exigir certificación de sitio. & Líder de Implantación & Certificación de sitio no emitida 8 semanas antes de la ola. \\
R-21 & Evitar & Validar los 25 supuestos y los umbrales con el cliente antes de la línea base. & Jefe de Proyecto & Supuestos sin validar en el mes 3. \\
R-22 & Mitigar & Asegurar acceso individualizado por proveedor y turnos de terminales compartidas. & Líder de Gestión del Cambio & Conteos incompletos en el piloto. \\
R-23 & Mitigar & Nombrar suplentes, documentar el conocimiento y adelantar la contratación. & Jefe de Proyecto & Vacante de un rol clave sin cubrir por más de 30 días. \\
R-24 & Mitigar & Ampliar el soporte reforzado y cerrar los defectos de severidad alta antes de abril de 2028. & Líder de Operación & Defectos altos abiertos a 4 semanas del Día de la Madre. \\
R-25 & Evitar & Cortar el último tramo antes del 15 de octubre, conciliar a diario y acordar un plan de retorno. & Líder de Datos y Migración & Algún tramo sin corte al 15 de octubre de 2028. \\
R-27 & Mitigar & Firmar acuerdos de cooperación y salida con documentación de interfaces y extracción de datos. & Jefe de Proyecto & Documentación sin entregar por un proveedor en el mes 4. \\
R-28 & Mitigar & Especificar la retención de grabaciones y acompañar la decisión del cliente de cambiar su política. & Arquitecto de Datos & Decisión del cliente sin tomar en el mes 6. \\
R-29 & Aceptar & Gestionar cualquier obligación nueva por el control de cambios y monitorear el procedimiento. & Jefe de Proyecto & Notificación de una obligación nueva. \\
R-30 & Evitar & Probar la creación de cada servicio Premium al inicio, antes de construir la plataforma base. & Arquitecto de Soluciones & Un servicio no creable en la prueba de recursos. \\
R-31 & Escalar & Resolver con el Comité Ejecutivo qué provee cada parte, cerrar la brecha de estándares y definir respaldo fuera de la comuna. & Responsable de Infraestructura & Informe de brecha de 2024 sin entregar en el mes 4. \\
R-32 & Mitigar & Confirmar ubicaciones y rutas con el inventario y acreditar el ExpressRoute antes del diseño final. & Responsable de Infraestructura & Inventario de sitios con cambios sobre el supuesto. \\
R-33 & Mitigar & Nivelar los recursos de los meses 13 a 15 y reforzar el equipo antes de la marcha blanca. & Jefe de Proyecto & Carga sobre 100 \% de un rol clave en los meses 13 a 15. \\
R-34 & Mitigar & Dimensionar las oleadas de cambio de etiquetas y evaluar las etiquetas electrónicas. & Líder de Operación & Error de precio sobre 3 \% en el piloto. \\
R-35 & Mitigar & Revisar los accesos periódicamente y revocar los de ex trabajadores y temporales. & Responsable de Seguridad & Cuenta de ex trabajador activa tras la revisión mensual. \\
R-36 & Mitigar & Incorporar transportistas y vendedores a los contratos de estado con calendario y medición. & Por definir con el sd-12 & Transportista sin adoptar los estados a 8 semanas del inicio. \\
R-37 & Evitar & Calendarizar las siete decisiones antes de la línea base y escalar las vencidas. & Jefe de Proyecto & Decisiones pendientes en el mes 3. \\
R-38 & Aceptar & Seleccionar pilotos y olas fuera del clúster de Santiago al inicio y revisar la concentración. & Líder de Implantación & Piloto con falla en una tienda del clúster de Santiago. \\
\end{longtable}

Las acciones se incorporan a la estructura de desglose del trabajo y los responsables son roles del equipo del proponente. La Tabla A.4 registra las oportunidades, que no se suman a las reservas.

Tabla A.4: Oportunidades. Fuente: Elaboración propia.

\begin{longtable}{p{0.074\linewidth}p{0.428\linewidth}p{0.112\linewidth}p{0.130\linewidth}p{0.186\linewidth}}
\toprule
\textbf{ID} & \textbf{Oportunidad} & \textbf{Probabilidad} & \textbf{Estrategia} & \textbf{Responsable} \\
\midrule
\endhead
\bottomrule
\endlastfoot
O-01 & Retirar antes alguna interfaz por archivo al cerrar el inventario de las 14 interfaces. & 0,40 & Explotar & Líder de Integración \\
O-02 & Adelantar la prueba de corte de 24 h si la filial fija los topes del cupo antes del mes 6. & 0,40 & Mejorar & Jefe de Proyecto \\
O-03 & Bajar la discrepancia del 12,4 \% antes del corte con el conteo cíclico en la carga inicial. & 0,40 & Explotar & Líder de Datos y Migración \\
O-04 & Detener el borrado a los 90 días si el Servicio de evidencia financiera entra antes del mes 16. & 0,20 & Mejorar & Arquitecto de Datos \\
O-05 & Ganar margen de observación si la cartera sin casos difíciles migra completa antes del mes 22. & 0,40 & Mejorar & Líder de Datos y Migración \\
O-06 & Dimensionar con medición si los datos del evento anual de 2026 calibran la capacidad. & 0,60 & Mejorar & Líder de Rendimiento y Calidad \\
O-07 & Adelantar la integración si los proveedores actuales entregan temprano la documentación de interfaces. & 0,20 & Compartir & Jefe de Proyecto \\
O-08 & Adelantar el despliegue si el cliente compra el hardware y contrata los enlaces antes de cada ola. & 0,20 & Compartir & Líder de Implantación \\
\end{longtable}

Ninguna de las ocho oportunidades se incluye en el cálculo de reservas del apartado 8.3.

\section*{Referencias}
\addcontentsline{toc}{section}{Referencias}

\noindent\hangindent=1.5em\hangafter=1 International Organization for Standardization. (2018). \textit{Risk management — Guidelines} (ISO 31000:2018).

\noindent\hangindent=1.5em\hangafter=1 Multitiendas Ancoa S.A. (2026). \textit{Bases técnicas. Caso 09: Cadena multitienda} (Licitación N.° TFEP-01/2026, versión 1.0).

\section*{Declaración de uso de IA}
\addcontentsline{toc}{section}{Declaración de uso de IA}

Este capítulo se redactó con asistencia de un modelo de inteligencia artificial (Claude, de Anthropic), a partir de decisiones tomadas por el equipo de Only Simple Solutions. Esas decisiones fijan el alcance del registro de riesgos, las escalas de probabilidad e impacto, la fusión de riesgos afines y los criterios de calificación.

El modelo generó con asistencia de IA el registro de riesgos, la calificación de cada riesgo y oportunidad, el análisis cuantitativo y las figuras del capítulo, así como el texto de los apartados y del anexo. Estos contenidos están pendientes de una revisión humana completa, por lo que las calificaciones, las estimaciones de semanas y horas, los umbrales y las fechas de disparador deben considerarse propuestas del equipo hasta que esa revisión se realice.

Las cifras y los hechos del Caso se tomaron de las bases técnicas citadas en Referencias. Las estimaciones y las escalas son decisiones del proponente y no provienen de esas bases.

\ossFinalPage

\end{document}
